\documentclass[12pt,a4paper]{article}
\usepackage[utf8]{inputenc}
\usepackage[indonesian]{babel}
\usepackage{amsmath}
\usepackage{amsfonts}
\usepackage{amssymb}
\usepackage{geometry}

\geometry{margin=2.5cm}

\title{Penyelesaian Titik Berat Bidang Polar Limacon Bagian Dalam}
\author{Ahmad Fakhrul Bawani}
\date{}

\begin{document}

\maketitle

\section*{Soal}

Diketahui persamaan kurva polar:
\[
r = 2\cos\theta + 1
\]
Kurva ini merupakan bentuk lima\c{c}on dengan loop dalam (\textit{inner loop}) karena konstanta lebih kecil dari koefisien trigonometri ($1 < 2$).

\subsection*{1. Rumus Utama Luas Polar}
Luas daerah ($A$) dalam koordinat polar yang disapu oleh jari-jari $r$ dari batas sudut $\alpha$ hingga $\beta$ ditentukan oleh integral:
\[
A = \frac{1}{2} \int_{\alpha}^{\beta} r^2 \, d\theta
\]

\subsection*{2. Menentukan Batas Sudut ($\alpha$ dan $\beta$)}
Loop dalam terbentuk ketika kurva melewati titik kutub asal, yaitu saat nilai jari-jari $r = 0$:
\[
2\cos\theta + 1 = 0 \implies \cos\theta = -\frac{1}{2}
\]
Pada rentang $[0, 2\pi]$, nilai sudut yang memenuhi persamaan tersebut adalah:
\[
\alpha = \frac{2\pi}{3} \quad \text{dan} \quad \beta = \frac{4\pi}{3}
\]

\subsection*{3. Formulasi dan Penyederhanaan Integral}
Substitusikan fungsi $r$ ke dalam rumus luas daerah:
\[
A = \frac{1}{2} \int_{\frac{2\pi}{3}}^{\frac{4\pi}{3}} (2\cos\theta + 1)^2 \, d\theta
\]
Jabarkan fungsi kuadrat di dalam integral:
\[
(2\cos\theta + 1)^2 = 4\cos^2\theta + 4\cos\theta + 1
\]
Gunakan identitas trigonometri sudut ganda $\cos^2\theta = \frac{1 + \cos(2\theta)}{2}$ untuk mereduksi pangkat suku pertama:
\[
4\left(\frac{1 + \cos(2\theta)}{2}\right) + 4\cos\theta + 1 = 3 + 4\cos\theta + 2\cos(2\theta)
\]

\subsection*{4. Evaluasi Integral}
Substitusikan kembali hasil penyederhanaan dan lakukan proses integrasi:
\[
A = \frac{1}{2} \int_{\frac{2\pi}{3}}^{\frac{4\pi}{3}} \left( 3 + 4\cos\theta + 2\cos(2\theta) \right) \, d\theta
\]
\[
A = \frac{1}{2} \left[ 3\theta + 4\sin\theta + \sin(2\theta) \right]_{\frac{2\pi}{3}}^{\frac{4\pi}{3}}
\]

Evaluasi nilai integral menggunakan batas atas dan batas bawah yang telah ditentukan:
\begin{itemize}
  \item Batas Atas ($\theta = \frac{4\pi}{3}$): $3\left(\frac{4\pi}{3}\right) + 4\sin\left(\frac{4\pi}{3}\right) + \sin\left(\frac{8\pi}{3}\right) = 4\pi - \frac{3}{2}\sqrt{3}$
  \item Batas Bawah ($\theta = \frac{2\pi}{3}$): $3\left(\frac{2\pi}{3}\right) + 4\sin\left(\frac{2\pi}{3}\right) + \sin\left(\frac{4\pi}{3}\right) = 2\pi + \frac{3}{2}\sqrt{3}$
\end{itemize}

Hitung selisih dari kedua hasil batas tersebut:
\[
A = \frac{1}{2} \left[ \left(4\pi - \frac{3}{2}\sqrt{3}\right) - \left(2\pi + \frac{3}{2}\sqrt{3}\right) \right]
\]
\[
A = \frac{1}{2} \left[ 2\pi - 3\sqrt{3} \right] = \pi - \frac{3}{2}\sqrt{3}
\]

\subsection*{Kesimpulan}
Luas daerah bagian dalam (\textit{inner loop}) dari lima\c{c}on $r = 2\cos\theta + 1$ adalah:
\[
A = \pi - \frac{3}{2}\sqrt{3} \approx 0.5435 \quad \text{satuan luas}
\]

\end{document}